\documentclass[10pt,letterpaper]{article}
\usepackage[margin=0.75in]{geometry}
\usepackage[T1]{fontenc}
\usepackage{lmodern,microtype,booktabs,tabularx,array,graphicx,amsmath,xcolor}
\usepackage{enumitem}
\usepackage[font=small,labelfont=bf]{caption}
\captionsetup{hypcap=false}
\usepackage[hidelinks]{hyperref}
\usepackage{xurl}
\definecolor{ink}{HTML}{20354C}
\usepackage{titlesec}
\titleformat{\section}{\large\bfseries\color{ink}}{\thesection}{0.6em}{}
\titleformat{\subsection}{\normalsize\bfseries\color{ink}}{\thesubsection}{0.6em}{}
\titlespacing*{\section}{0pt}{11pt}{5pt}
\titlespacing*{\subsection}{0pt}{8pt}{3pt}
\setlength{\parindent}{0pt}
\setlength{\parskip}{4pt}
\setlength{\emergencystretch}{2em}
\setlist[itemize]{leftmargin=1.25em,itemsep=2pt,topsep=3pt}
\hypersetup{pdftitle={Real-Time Gym Form Correction: Midpoint Report},pdfauthor={Ellis Charles and Kevin Xue}}
% Generated assets live beside the packaged source, or under data/reports when
% compiling this editable source from docs/report or the repository root.
\newcommand{\reportassets}{}
\IfFileExists{tables/metrics_macros.tex}{}{%
  \IfFileExists{../../data/reports/midpoint/tables/metrics_macros.tex}{%
    \renewcommand{\reportassets}{../../data/reports/midpoint/}%
  }{%
    \IfFileExists{data/reports/midpoint/tables/metrics_macros.tex}{%
      \renewcommand{\reportassets}{data/reports/midpoint/}%
    }{%
      \PackageError{midpoint}{Generated report assets are missing}{Run make figures from the repository root, then compile again.}%
    }%
  }%
}
\graphicspath{{\reportassets figures/}}
\input{\reportassets tables/metrics_macros.tex}
\begin{document}
\begin{center}
{\LARGE\bfseries\color{ink} Real-Time Gym Form Correction}\par
\vspace{5pt}
{\large CIS 5810 Midpoint Report --- Group 53}\par
\vspace{5pt}
Ellis Charles: \texttt{chaelli@engineering.upenn.edu}\\
Kevin Xue: \texttt{kevinxue@engineering.upenn.edu}\\
\vspace{3pt}
November 22, 2026
\end{center}

\section*{Summary}
We are building a webcam-based exercise assistant that counts squats and gives
feedback on squat depth. We tested five versions on six development clips
containing 17 manually counted repetitions. Smoothing the knee-angle signal
reduced mean absolute counting error (MAE) from \BaseMAE{} to \SmoothMAE{} reps
per clip. We then tried using a reference squat to judge depth and added checks
on movement duration and tracking quality. Both versions had a counting MAE of
\RefMAE{}, but their depth counts were less accurate than the baseline's.
Our latest version combines 2D knee-angle tracking for counting with
reference-based 3D features for depth feedback. It achieved a counting MAE of
\HybridMAE{}, although its depth-count MAE of \HybridDepthMAE{} was still worse
than the baseline's \BaseDepthMAE{}. It also marked depth as unclear for two of
its 13 counted reps. We developed this version after reviewing the same six
clips, so we still need to test it on new videos.

\section{Goals, audience, and scope}
Our goal is to help people who exercise without a trainer. The planned app
shows a pose skeleton, a rep count, and depth feedback over a camera feed.
For the midpoint, we focus on counting completed squats and counting how many
reach acceptable depth. We define acceptable depth as the hip crease reaching
the top of the knee. This evaluation does not cover injury risk, overall squat
technique, or recognition of other exercises.

\section{System and pipeline}
We use the pretrained MediaPipe Pose Landmarker Lite model to estimate 2D image
landmarks and 3D world landmarks. We do not retrain the pose model.
BlazePose describes the approach behind this type of pose tracking for fitness
applications~\cite{bazarevsky2020,mediapipe}. Our work focuses on processing the
landmarks, tracking each squat, calibrating depth feedback, and evaluating the
results.
\begin{center}
\includegraphics[width=\linewidth]{pipeline.png}
\captionof{figure}{All five methods use the same saved pose data.
The manual counts are used to evaluate predictions.}
\end{center}
We save the pose results so every method uses the same inputs. These files
include frame timestamps, missing observations, landmarks, visibility scores,
computed features, and processing times. We also save hashes to identify the
video, model, and extraction code used. The repository separates reusable code,
scripts, data, and generated results.

\subsection*{Baseline}
The baseline uses the knee-angle state machine from our original proposal.
It tracks the more visible leg, starts a rep when the knee angle falls below
$150^\circ$, and counts the rep when the angle returns above $160^\circ$.
It marks depth as acceptable if the smallest 2D knee angle during the rep is
at most $95^\circ$. We use these thresholds as starting values; they have not
been validated as anatomical standards.

\clearpage
\section{Data and evaluation}
We use videos from Khorasani's \emph{Gym Members Exercise Dataset} on
Kaggle~\cite{khorasani2024}. Our local archive contains 29 squat clips totaling
about 297 seconds. We reviewed sample frames to check for cropping, split-screen
layouts, edits, and footage we could use. We chose one reference clip and six
development clips, keeping the original video files. The six evaluation clips
contain 17 completed reps, with 15 labeled as reaching acceptable depth.
The reference clip, clip 11, is used for calibration and is not scored.

\begin{center}
\input{\reportassets tables/dataset.tex}
\captionof{table}{Clips used in the experiment. Depth coverage is the percentage
of frames containing all three features needed by the reference method.
It does not measure depth accuracy. Durations come from the decoded timestamps.}
\end{center}

For each evaluation clip, we recorded the total number of completed reps and
the number that reached acceptable depth. We excluded partial reps at the
start or end of a video. For example, clip 4 has four completed reps and three
acceptable-depth reps, with the annotation ``first one bad.'' Both team members
reviewed the depth rule and a subset of the counts, but we did not measure
agreement between annotators. Depth is also harder to judge from the front
than from the side.

In clip 11, we marked the start, bottom, and return of one acceptable squat
at 0.539, 2.523, and 4.830 seconds. We use only this interval to calculate the
reference standing and bottom positions. The source groups are
\texttt{gym\_b} for calibration and \texttt{gym\_a}, \texttt{tutorial\_c}, and
\texttt{tutorial\_d} for development. The evaluation script checks for duplicate
video files and prevents a source group from appearing in both splits.
However, the participant field contains the placeholder \texttt{video}, so we
cannot confirm that the splits contain different people.

For the $N=6$ evaluation clips, we report mean absolute counting error and the
fraction of clips whose predicted count matches the manual count exactly:
\[
\mathrm{MAE}=\frac{1}{N}\sum_{i=1}^{N}|\widehat n_i-n_i|,
\qquad
\mathrm{Exact}=\frac{1}{N}\sum_{i=1}^{N}\mathbf{1}[\widehat n_i=n_i].
\]
We calculate depth-count MAE the same way, using the number of acceptable-depth
reps. Since our labels are totals for each video, we cannot report per-rep
precision, recall, F1, or depth-classification accuracy. We also have no clips
with zero squats, so we have not tested whether standing or walking produces
false counts.

The first four methods use the same saved pose data and fixed settings.
We developed the hybrid method after reviewing their errors. It keeps the
0.6-second minimum duration and 0.25-second gap limit, but its design still
reflects what we learned from these clips. All six scored clips are marked
\texttt{dev}; we have not evaluated on a separate test set. The 3D coordinates
are estimates from a single camera, rather than motion-capture
measurements~\cite{mediapipe}.

\clearpage
\section{Baseline and first three iterations}
\begin{center}
\input{\reportassets tables/results.tex}
\captionof{table}{Results on the same six development clips. Lower MAE is
better. The reference and temporal methods produce identical predictions.
The hybrid method, described on the next page, was developed after reviewing
these results. ``Unclear'' counts reps whose depth the method could not judge.}
\end{center}

\subsection{Iteration 1: smoothing the 2D signal}
The baseline overcounts all three gym clips filmed at an angle. In clip 5,
for example, it counts 10 reps instead of four. We added exponential smoothing
with a time constant of 0.08 seconds, using only current and previous frames.
We also cancel an unfinished rep if tracking is missing for more than
0.25 seconds. Counting MAE drops from \BaseMAE{} to \SmoothMAE{}, a
\CountReduction{}\% reduction. In clip 5, the count drops from 10 to five.
The number of exact matches stays at 1/6, while depth-count MAE increases from
\BaseDepthMAE{} to \SmoothDepthMAE{}. Smoothing reduces extra counts, but depth
judgments still depend on the camera angle.

\subsection{Iteration 2: using a reference squat for depth}
We replaced the 2D knee-angle signal with hip height relative to torso length
and both 3D knee angles. For each feature $q_j$, we use standing value $s_j$ and
bottom value $b_j$ from the reference squat to calculate
$r_j=(s_j-q_j)/(s_j-b_j)$. We combine them into a depth signal:
$d=0.5r_{\mathrm{hip}}+0.5\min(r_{\mathrm{left}},r_{\mathrm{right}})$.
A rep starts at $d\geq0.25$ and ends after returning to $d\leq0.12$.
Its peak must reach at least 0.8 for depth to be accepted. Counting MAE stays at
\RefMAE{}, and exact matches increase to 2/6. Clip 18 improves from zero to one
counted rep, but clip 26 drops from three to zero. Depth-count MAE increases to
\RefDepthMAE{}. These results suggest that one reference squat does not transfer
well across the people and camera views in our clips.

\subsection{Iteration 3: checking movement over time}
We added a minimum rep duration of 0.6 seconds, a check that the person rises
from the bottom, and a requirement to stay at depth for at least 0.08 seconds.
We also mark depth as unclear when the required features are available for
less than 80\% of the rep. These checks aim to reject brief false movements and
avoid giving depth feedback when tracking is incomplete. However, every rep
counted by the reference method already passes these checks. Counts and depth
labels therefore stay the same on our six clips. The automated tests check
that the rules work as intended, but we have not yet shown that they improve
results on real videos.

\begin{center}
\includegraphics[width=\linewidth]{method_comparison.png}
\captionof{figure}{The updated methods reduce counting error, but all have
higher depth-count error than the baseline.}
\end{center}

\clearpage
\section{Iteration 4: separating rep counting from depth feedback}
The reference method can miss an entire rep when it loses track of one leg.
To address this, our hybrid method counts reps using the smoothed 2D knee
angle and uses the reference-based 3D features only to judge depth.
Before starting a rep, it must first observe the person standing, including
after a tracking gap. A completed rep must return to standing, last
0.6--15 seconds, and have usable 2D tracking for at least 80\% of its duration.
A tracking gap longer than 0.25 seconds cancels the unfinished rep.

The method can count a completed rep even when it cannot judge its depth.
It marks depth as \emph{unclear} if too many 3D features are missing, if there
is a gap longer than 0.25 seconds, or if the apparent depth peak is too brief.
We do not fill in missing features. Only reps explicitly marked as acceptable
contribute to the accepted-depth count; an unclear label does not mean that
the squat was shallow.

\begin{center}
\input{\reportassets tables/counts.tex}
\captionof{table}{Predicted rep counts for each clip. The hybrid counts all
three reps in clip 26, but misses clip 18 and undercounts the three angled gym clips.}
\end{center}

The hybrid counts 3, 3, 3, 1, 0, and 3 reps, compared with manual counts of
4, 4, 4, 1, 1, and 3. Its counting MAE is \HybridMAE{}, compared with
\SmoothMAE{} for the smoothing method, and it matches two of the six clips
exactly. Two of its 13 counted reps have unclear depth, leaving 84.6\% with
an acceptable or shallow label. Depth-count MAE is \HybridDepthMAE{}, so better
rep counting has not yet led to better depth judgments.
In an earlier trial, adding only the duration check gave a counting MAE of
0.50. The full hybrid scores 0.67 because its standing, gap, and coverage checks
also reject interrupted movements. The tables report the full implementation.

\begin{center}
\includegraphics[width=\linewidth]{hybrid_case.png}
\captionof{figure}{Replaying clip 26 through the demo counts three reps.
Two have acceptable depth and one has unclear depth. This checks that the
tracker works in the video demo; it does not measure live-camera latency.}
\end{center}

\clearpage
\section{Errors and processing speed}
\begin{center}
\includegraphics[width=\linewidth]{failure_traces.png}
\captionof{figure}{Top: the 2D knee angle repeatedly crosses the counting
thresholds in clip 5. Bottom: clip 26 is missing features that the reference
method needs from both legs. Shaded areas show missing features.}
\end{center}

\textbf{Repeated threshold crossings.} Eight of the baseline's ten counted
reps in clip 5 last less than 0.6 seconds. Together with the plot, this suggests
that noisy angle estimates cause extra counts. We would need labels for each
rep to identify every false count. Switching the tracked leg between frames,
background distractions, and blocked body parts may also contribute.

\textbf{Missing features.} Clip 26 has usable depth features in only 56.25\%
of frames. The longest gap between valid observations is 2.04 seconds, well
above our 0.25-second limit. The reference tracker is not ready to start a rep
in 88.0\% of frames and counts no reps, while the 2D baseline counts all three.
Clip 17 has similar problems: 40\% depth coverage and a 1.72-second gap between
valid observations.

\textbf{Reference calibration.} The largest raw depth values in clips 2 and 4
are 0.45 and 0.65, below the required 0.8. However, the manual labels mark four
and three reps as reaching acceptable depth. The reference-based signal
therefore disagrees with the labels on these clips. We need to review more examples to understand this mismatch.

\textbf{Processing speed.} The saved median pose and feature extraction times
range from \ExtractMin{} to \ExtractMax{} ms per frame across clips on our local
macOS arm64 machine. Median counter times from a new run are about
\CounterMin{}--\CounterMax{} ms per frame, although these very small measurements
are sensitive to timer resolution. These times exclude video decoding,
display, and camera buffering. We still need to measure the delay for the
complete webcam demo.

\clearpage
\section{Next steps, timeline, and responsibilities}
We have implemented and compared five squat-counting methods. The next phase
will focus on depth errors and evaluation on new clips. We have not yet compared
against MoveNet or tested recognition of other exercises. The demo uses the
hybrid by default, with the original temporal method still available. Video
replays count three reps in clip 26 and one in clip 17 without reference
calibration. Without a reference, depth is marked unclear. We also need to
record a backup presentation video.

Our next steps focus on the errors we observed and on testing new clips:
\begin{itemize}
\item Improve depth estimates when some landmarks are missing, and make person
and leg selection more stable. The hybrid keeps more rep counts, but still
returns unclear or incorrect depth labels.
\item Try calibration with several acceptable squats or camera views, then test
whether it works on source groups we have not used for development.
\item Add clips with brief movements, tracking noise, standing, and walking.
Use them to check whether the timing rules reject false counts without
removing valid reps.
\end{itemize}

\begin{center}
\small
\begin{tabularx}{\linewidth}{@{}p{0.19\linewidth}Xp{0.19\linewidth}@{}}
\toprule
Target & Milestone & Lead / support\\
\midrule
Nov. 22 & Finish the five-method comparison and revised midpoint report & Both\\
Nov. 23--25 & Review missing landmarks, improve unclear-depth handling, add partial-rep and no-squat clips & Ellis / Kevin\\
Nov. 26--29 & Test calibration with multiple reference squats and timing checks on new development clips & Kevin / Ellis\\
Nov. 30--Dec. 3 & Fix settings, label new source groups, evaluate, and measure webcam latency & Both\\
Dec. 4--6 & Rehearse the five-minute presentation, record a backup, and attend Demo Day on Dec. 6 & Both\\
Dec. 7--10 & Analyze the final experiments and finish the report and supporting materials by Dec. 10 & Both\\
\bottomrule
\end{tabularx}
\end{center}

Following our proposal, Ellis leads pose-estimation integration and video
processing, while Kevin leads squat-state tracking and feedback. Both members
help review annotations, connect the system components, interpret results,
and prepare the presentation. The table shows our planned work for the
rest of the project.

\subsection*{Reproducing the results and crediting sources}
The source and results package includes the annotation CSVs, settings,
reference values, per-frame predictions, file and code hashes, and comparison
tables. Running \texttt{make report} evaluates all five methods and builds the
figures, LaTeX report, and supporting ZIP. Running \texttt{make check} runs the
automated tests. All 57 tests pass, which checks software behavior rather than
accuracy on additional videos. We credit MediaPipe/BlazePose and the video
dataset below. Our earlier MM-Fit data review and phase-template prototype
are not part of the results reported here.


\begingroup
\small
\bibliographystyle{plain}
\IfFileExists{references.bib}{\bibliography{references}}{\bibliography{docs/report/references}}
\endgroup
\end{document}
